% TOP_PROBLEM: {"id": 30006769, "title": "Berman-Hartmanis Isomorphism Conjecture", "category_id": 15, "category": {"id": 15, "name": "computer_science", "display_name": "Computer Science", "description": "Computational complexity, algorithms, and theoretical CS.", "slug": "computer-science", "order_index": 15, "created_at": "2026-07-31T15:26:25.671Z"}, "status": "open"}
% FIRST_POSTED: 2026-09-27
% !TeX program = lualatex
% Compile twice: lualatex 30006769.tex
% All references and prose are embedded; no BibTeX database or external inputs.
\documentclass[11pt,a4paper]{article}
\usepackage[margin=24mm,headheight=14pt]{geometry}
\usepackage{fontspec}
\setmainfont{Latin Modern Roman}
\setsansfont{Latin Modern Sans}
\setmonofont{Latin Modern Mono}
\usepackage{amsmath,amssymb,mathtools}
\usepackage{unicode-math}
\setmathfont{Latin Modern Math}
\usepackage{microtype}
\usepackage{enumitem,xurl,needspace}
\usepackage[dvipsnames]{xcolor}
\usepackage{hyperref}
\hypersetup{unicode=true,colorlinks=true,linkcolor=MidnightBlue,urlcolor=MidnightBlue,
 pdftitle={500 Mathematical Problem Statements},pdfauthor={Source-annotated compilation},bookmarksnumbered=true}
\usepackage{bookmark}
\usepackage{tocloft}
\setlength{\cftsecnumwidth}{3em}
\cftsetpnumwidth{2.5em}
\cftsetrmarg{4em}
\usepackage{titlesec}
\titleformat{\section}{\Large\bfseries\raggedright}{\thesection}{0.7em}{}
\usepackage{fancyhdr}
\pagestyle{fancy}\fancyhf{}
\fancyhead[L]{\small\sffamily Mathematical problem statements}
\fancyhead[R]{\small Source edition 22}
\fancyfoot[C]{\thepage}
\setlength{\parindent}{0pt}
\setlength{\parskip}{5pt plus 1pt}
\setlength{\emergencystretch}{3em}
\setcounter{tocdepth}{1}
\setcounter{secnumdepth}{1}
\setlist[itemize]{leftmargin=3em,itemsep=5pt,topsep=4pt,parsep=0pt}
\allowdisplaybreaks
\newcommand{\N}{\mathbb{N}}
\newcommand{\Z}{\mathbb{Z}}
\newcommand{\Q}{\mathbb{Q}}
\newcommand{\R}{\mathbb{R}}
\newcommand{\C}{\mathbb{C}}
\newcommand{\F}{\mathbb{F}}
\newcommand{\A}{\mathbb{A}}
\newcommand{\PP}{\mathbb P}
\newcommand{\E}{\mathbb{E}}
\newcommand{\eps}{\varepsilon}
\newcommand{\dd}{\,\mathrm d}
\newcommand{\sref}[2]{\hyperlink{source:#1:#2}{[S#2]}}
\newcommand{\eref}[2]{\hyperlink{extra:#1:#2}{[E#2]}}
% Turn the next switch off for a shorter reading copy. The quotations preserve
% every exactTarget field, including qualifications omitted from a short summary.
\newif\ifshowcatalogue
\showcataloguetrue
\newcommand{\cataloguescope}[1]{\ifshowcatalogue
 \paragraph{Exact scope in the supplied catalogue.}
 \begin{quote}\small #1\end{quote}\fi}

% Standalone notebook wrapper; original problem section retained verbatim.
\numberwithin{equation}{section}
\hypersetup{pdftitle={Research attempt -- problem 30006769},
 pdfauthor={Research notebook; source compilation retained}}
\fancyhead[L]{\small\sffamily Research notebook}
\fancyhead[R]{\small\sffamily Problem 30006769 | 27 September 2026}
\begin{document}
\setcounter{section}{0}
% ================================================================
% problemId: 30006769
% Canonical problem ID: 30006769
% exactTarget SHA-256: 81865b4ee16fbd06fc1b3d28e9643b56c901af28876edece420e3847ea10003f
\clearpage
\section*{\texorpdfstring{Berman-Hartmanis Isomorphism Conjecture}{Berman-Hartmanis Isomorphism Conjecture}}\label{30006769}
\subsection{Definitions and mathematical statement}
A language $A\subseteq\{0,1\}^*$ is NP-complete under polynomial-time many-one reductions if $A\in\mathrm{NP}$ and every NP language $L$ has a polynomial-time map $f$ with $x\in L\iff f(x)\in A$. A polynomial-time isomorphism between $A$ and $B$ is a bijection $h:\{0,1\}^*\to\{0,1\}^*$ satisfying
\[
 x\in A\iff h(x)\in B,
 \qquad h,h^{-1}\text{ polynomial-time computable}.
\]
The conjecture asserts such an $h$ for every pair of NP-complete languages. Reductions in both directions alone do not supply a bijection or an efficiently computable inverse.
\par\smallskip\textit{Formulation and concept references:} \sref{30006769}{1}.
\cataloguescope{Prove or disprove that for every two languages A and B complete for NP under polynomial-time many-one reductions, there is a bijection h on binary strings such that h(A)=B and both h and its inverse are computable in polynomial time.}
\subsection{Short English statement}
Are all NP-complete languages the same up to a reversible, polynomial-time change of their binary encodings?
% BEGIN RESEARCH ADDITION ATTEMPT-175
\subsection{Research attempt}

\paragraph{Current frontier.}
The Hartmanis problem survey discusses the isomorphism conjecture, padding,
scrambling, and restricted reductions separately \sref{30006769}{1},
Sections 5.1--5.3. In particular, results for completeness under weaker
reduction notions do not establish the statement for all polynomial-time
many-one complete sets. Its discussion also records relativized obstacles.
The present attempt isolates a sufficient effective Cantor--Bernstein
mechanism and checks exactly which features arbitrary Karp reductions
lack. It does not assume that injectivity supplies efficient inversion.

\paragraph{Attack route.}
One route tries to pad every complete language reversibly; another tries to
convert reductions in both directions directly into a bijection. A purely
set-theoretic Cantor--Bernstein theorem ignores the cost of recognizing the
components used in its proof. We pursue a version in which that recognition
is forced to terminate quickly by a decreasing integer measure: string
length. The remaining task is then an explicit upgrade theorem for
reductions, rather than an unspecified ``back-and-forth'' construction.

\paragraph{Attempt: effective Cantor--Bernstein under explicit hypotheses.}
Let $A,B\subseteq\{0,1\}^*$. Suppose there are injections $f,g$ with
\[
 x\in A\iff f(x)\in B,\qquad y\in B\iff g(y)\in A.
\]
Assume $f,g$ are polynomial-time computable, satisfy
\begin{equation}\label{eq175-length}
              |f(x)|>|x|,\qquad |g(y)|>|y|,
\end{equation}
and have polynomial-time partial inverses: given a string, an algorithm
recognizes membership in the range and either returns the unique preimage
or reports that none exists. Then $A$ and $B$ are polynomial-time
isomorphic in the exact sense of the question.

\textit{Proof.} Make two disjoint copies $X,Y$ of the set of binary strings.
Draw an edge from $x\in X$ to $f(x)\in Y$ and from $y\in Y$ to
$g(y)\in X$. Every vertex has exactly one outgoing edge and at most one
incoming edge. By \eqref{eq175-length}, following edges backwards strictly
decreases string length, so every component has a unique initial vertex,
called its root. There are no cycles or doubly infinite chains. Each
component is a one-sided alternating ray starting in $X$ or in $Y$.

The root side of a vertex is polynomial-time decidable. Alternate the
range tests for $f$ and $g$ as appropriate; each successful inversion
shortens the string by at least one. There are at most the original length
plus one tests, each on strings no longer than the original. Stop at the
first failed range test. Polynomial running time follows with a fixed
polynomial depending on the two given inverse algorithms.

Define
\begin{equation}\label{eq175-h}
 h(x)=\begin{cases}
       f(x),&\text{the root of its component is in }X,\\
       g^{-1}(x),&\text{the root is in }Y.
      \end{cases}
\end{equation}
In the second case every $X$-vertex has a predecessor in $Y$, so the
inverse is defined. On a ray rooted in $X$, this matches successive
$X$-vertices with their $f$-successors. On a ray rooted in $Y$, it matches
each $X$-vertex with its $g$-predecessor. In either case it is a bijection
between the two sets of vertices of the ray. The inverse is
$f^{-1}(y)$ on an $X$-rooted ray and $g(y)$ on a $Y$-rooted ray.
Both maps are polynomial-time by the root algorithm and the given maps.
All their paired edges preserve the relevant membership condition, proving
$h(A)=B$. This proves the lemma. \hfill$\square$

The output length is also controlled: it is either the output of one
polynomial-time map or a shorter predecessor. No infinite component search
or hidden oracle has been used.

\paragraph{Attempt: what padding would have to supply.}
The lemma would prove the full conjecture if every pair of NP-complete
languages admitted mutual reductions satisfying its three extra properties:
injectivity, effective range inversion, and strict length increase. The
existence of ordinary many-one reductions gives none of those properties
for free. Two different inputs may have the same image, and their membership
labels alone provide no reversible encoding of the discarded information.

For a paddable target, one might try to replace a reduction $r(x)$ by an
encoding $\operatorname{pad}(r(x),x)$. This works only if padding preserves
membership and the pair can be decoded efficiently from the \emph{target
language's own strings}. It is not enough to form the pair in an external
cylinder language and then apply another arbitrary reduction: that last
reduction may destroy the very injectivity and decoding that were gained.
The role of paddability and scrambling in this obstruction is discussed in
\sref{30006769}{1}, Sections 5.1--5.2.

Likewise, an injection can be difficult to invert on its range. In that
case the predecessor tests in the proof are not algorithms of known
polynomial complexity. If length increase is dropped, backwards paths
may not terminate, and even a finite component can contain a cycle.
An effective treatment of those components might still be possible in
special cases, but it requires a new argument; the proof above does not
cover them.

As a mechanical test of the pairing construction, take $f(x)=0x$ and
$g(y)=1y$. Range inversion deletes the appropriate leading bit. The
accompanying script constructs \eqref{eq175-h} and its inverse and checks
both compositions on all strings up to a specified finite length. The
all-string proof remains the ray argument, not this finite test. These
prefix maps test the combinatorics, not arbitrary NP-complete languages.

\paragraph{Attempt: a dependency warning proved directly.}
An affirmative answer to the full question would imply
$\mathrm P\ne\mathrm{NP}$. Suppose instead that $\mathrm P=\mathrm{NP}$.
Then every nonempty proper language in $\mathrm P$ is NP-complete: given
any NP language, decide it in polynomial time and output one fixed yes
string or one fixed no string of the target. In particular
\[
 A=\{0\},\qquad B=\{0x:x\in\{0,1\}^*\}
\]
would both be NP-complete. But no bijection of all strings can carry the
singleton $A$ onto the infinite set $B$. Thus the universal isomorphism
assertion fails under $\mathrm P=\mathrm{NP}$.

This uses the exact definition in the question, which does not exclude
finite complete languages by an extra convention. It confirms the
warning in \sref{30006769}{1}. A theorem upgrading all complete sets to the
reversible embeddings above is therefore not a modest coding exercise:
it would already imply a major complexity separation. Conversely, a
failure to upgrade a particular chosen pair of reductions would not
disprove the conjecture, which permits entirely different isomorphisms.

\paragraph{Stress test.}
The bijection acts on \emph{all} binary strings, including the empty string,
and not just on the yes instances. Roots of length zero cause no problem:
no shorter predecessor can exist. Membership equivalences are used in
both directions along each edge, so the inverse also preserves the no
instances. Fixed polynomials bound the running times for all inputs;
nonuniform inverse circuits are not substituted for uniform algorithms.

Set-theoretic reducibility and polynomial-time invertibility remain distinct.
The source's oracle and scrambling discussions prevent treating this proof
as a general relativizing solution. No one-way function is assumed to exist,
and no unproved cryptographic hypothesis is smuggled into the status of
the upgrade step.

\paragraph{Outcome.}
\textbf{Outcome: Conditional reduction.}

The attempt gives a complete constructive proof of an efficient
Cantor--Bernstein lemma with explicit length and inversion conditions.
The missing input is a theorem supplying those conditions for arbitrary
NP-complete languages. The lemma is an exposition and audit of a standard
kind of isomorphism mechanism, not a claim of literature priority. Its
application to the universal conjecture remains conditional, and its
consequence for $\mathrm P$ versus $\mathrm{NP}$ is explicitly verified.
% END RESEARCH ADDITION ATTEMPT-175
\subsection{Sources}
\begin{itemize}\raggedright
\item[\textnormal{[S1]}]\hypertarget{source:30006769:1}{} Allender, Cai, Fortnow, Gasarch, Immerman, Kurtz, Royer, and Williams, Open Problems by or Inspired by Juris Hartmanis, ACM SIGACT News 53(4), 2022, section 5.3, printed p. 7.
\url{https://www.cs.umd.edu/~gasarch/open/juris.pdf}.
{\small\textit{Catalogue formulation source.}}
\end{itemize}

\end{document}
